\documentclass[10pt,a4paper]{article}
\usepackage{amsmath,amssymb,amsthm,graphicx}
\usepackage{amsmath,amssymb,amsthm,hyperref,xfrac}
\providecommand{\ie}{i.e.\ }
\providecommand{\skpt}{\smallskip}
\usepackage[scr=rsfs,cal=euler]{mathalfa}
\multlinegap0pt

\hyphenation{begin decom-po-si-tion sequence}
\newenvironment{enumeratea}
{\bgroup\def\theenumi{\alph{enumi}}\begin{enumerate}}
{\end{enumerate}\egroup}

\newcommand{\ba}{{\boldsymbol a}}
\newcommand{\bb}{{\boldsymbol b}}
\newcommand{\bd}{{\boldsymbol d}}
\newcommand{\bu}{{\boldsymbol u}}
\newcommand{\bx}{{\boldsymbol x}}
\newcommand{\bmy}{{\boldsymbol y}}

\newcommand{\bD}{{\boldsymbol D}}
\newcommand{\bL}{{\boldsymbol L}}
\newcommand{\bP}{{\boldsymbol P}}
\newcommand{\bV}{{\boldsymbol V}}
\newcommand{\bX}{{\boldsymbol X}}
\newcommand{\bZ}{{\boldsymbol Z}}

\newcommand{\Psfrac}[2]{(\sfrac{#1}{#2})}
\newcommand{\psfrac}[2]{\sfrac{(#1)}{#2}}
\newcommand{\bigO}[1]{O(#1)}

\usepackage{mathtools}
\usepackage{bbm}

\newcommand{\NN}{\mathbb{N}}
\newcommand{\ZZ}{\mathbb{Z}}
\newcommand{\RR}{\mathbb{R}}
\newcommand{\CC}{\mathbb{C}}
\newcommand{\PP}{\mathbb{P}}
\newcommand{\EE}{\mathbb{E}}

\newcommand{\boldalpha}{\boldsymbol{\alpha}}
\newcommand{\bolddelta}{\boldsymbol{\delta}}
\newcommand{\boldrho}{\boldsymbol{\rho}}
\newcommand{\dee}{\,\mathrm{d}}
\newcommand{\one}{\mathbbm{1}}
\newcommand{\BC}{\mathcal{B}}
\newcommand{\EC}{\mathcal{E}}
\newcommand{\RS}{\mathscr{R}}
\newcommand{\prob}[1]{\PP[ #1 ]}
\newcommand{\bigoo}[1]{O\!\left(#1 \right)}
\newcommand{\abs}[1]{\left\lvert #1 \right\rvert}
\newcommand{\floor}[1]{\left\lfloor #1 \right\rfloor}
\newcommand{\riemint}[4]{\int_{#1}^{#2}\! #3 \, \mathrm{d}#4}
\newcommand{\maxnorm}[1]{\left\| #1 \right\|_{\infty}}

\renewcommand{\Re}{\textup{Re}}
\renewcommand{\Im}{\textup{Im}}

\newcommand{\dtv}{\mathrm{d}_{\mathrm{TV}}}
\newcommand{\dw}{\mathrm{d}_{\mathrm W}}

\let\epsilon\varepsilon

\newcommand{\pextra}{P_{\mathtt{extra}}}
\newcommand{\parratia}{P_{\mathtt{Arratia}}}

\newcommand{\simplex}{\Delta^{k-1}}

\newcommand{\bg}{\big}
\newcommand{\bgg}{\Big}
\newcommand{\bggg}{\bigg}

\newcommand{\Rnt}{R_{\mathrm{NT}}}
\newcommand{\Rprob}{R_{\mathrm{PD}}}
\newcommand{\ECnt}{\EC_{\mathrm{NT}}}
\newcommand{\ECprob}{\EC_{\mathrm{PD}}}

\newcommand{\Bnt}{\mathcal{B}_{\mathrm{NT}}}
\newcommand{\Bprob}{\mathcal{B}_{\mathrm{PD}}}

\DeclareMathOperator{\vard}{Var}
\DeclareMathOperator{\Dir}{Dir}

\theoremstyle{plain}
\newtheorem{cor}{Corollary}[section]
\newtheorem{lem}[cor]{Lemma}
\newtheorem{prop}[cor]{Proposition}
\newtheorem{conj}[cor]{Conjecture}
\newtheorem{theo}{Theorem}

\theoremstyle{definition}
\newtheorem{example}{Example}
\newtheorem{deff}[cor]{Definition}
\newtheorem{rem}{Remark}[section]
\newtheorem*{rems}{Remark}
\newtheorem*{remss}{Remarks}

\makeatletter
\ams@newcommand{\multiint}[1]{\DOTSI\protect\MultiIntegral{#1}}
\renewcommand{\MultiIntegral}[1]{%
\edef\ints@c{\noexpand\intop
\ifnum#1=\z@\noexpand\intdots@\else\noexpand\intkern@\fi
\replicate{#1-2}{\noexpand\intop\noexpand\intkern@}%
\noexpand\intop
\noexpand\ilimits@
}%
\futurelet\@let@token\ints@a
}
\makeatother
\ExplSyntaxOn
\cs_new:Npn \replicate #1 #2 { \prg_replicate:nn { #1 } { #2 } }
\ExplSyntaxOff
\title{The first coupling event fails with probability of order one over the logarithm}
\author{Tony Haddad, Dimitris Koukoulopoulos}
\date{doi:10.5802/jep.317; CC BY 4.0}
\begin{document}
\maketitle
\noindent\textit{Editorial note.} The counted statement is the author's labelled lem environment at source0.tex lines 886-889; its proof environment opens at line 891, the line after the blank separator, so the pairing is by adjacency and the retained proof range 891-961 is the whole environment.\par\medskip
\noindent\textit{Editorial note.} Four environments are retained as context because the counted proof depends on them: eq:mgf theta infty (source lines 728-731), the scale-invariant spacing lemma (source lines 695-698), prop:icexponential (source lines 814-817) and the Mecke equation prel:mecke (source lines 1845-1852).\par\medskip
\noindent\textit{Editorial note.} This article carries no label-number dump and the retained counter values are left empty for the whole source, so the printed numbering of the card is the wrapper's default; the identity of the counted item is its author label lem:prob e1c.\par\medskip
\noindent\textit{Editorial scope.} One original result of the pinned publisher article, counted once: the estimate of order one over the logarithm for the failure of the first coupling event (author label lem:prob e1c).\par\medskip
\noindent\textit{Allowed standard prerequisites (editorial).}\begin{itemize}
\item The first coupling event. The event E1 controls the total mass of the point process and fails when some level exceeds the logarithm of x.. It is the event whose complement is estimated.
\item Moments of the limiting variable. The exponential moment of the limiting variable is bounded uniformly, hence the probability that it exceeds the iterated logarithm is of the order one over the logarithm.. It removes the region where the level exceeds the iterated logarithm.
\item Exponential law of the truncated weighted sum. The truncated weighted sum is exponentially distributed with parameter equal to the level.. It bounds the tail probability appearing after the expectation is computed.
\end{itemize}\medskip
\begin{equation}
\label{eq:mgf theta infty}
\EE[e^{\alpha \Theta_\infty}] = \exp\!\bggg(\!\int_{0}^\infty \frac{e^{\alpha r(y)} - 1}{y}\, \mathrm dy\bggg).
\end{equation}

\begin{lem}[Scale-invariant spacing lemma]
\label{lem:scale-invariant-spacing-lemma}
The point process $\{S_i : i \in \ZZ\}$ is a Poisson process on the positive real line with intensity measure $\sfrac{\mathrm ds}{s}$.
\end{lem}

\begin{prop}[Arratia, \cite{Arratia02}]
\label{prop:icexponential}
For any fixed $t > 0$, the random variable $I_t$ follows an exponential distribution of parameter $t$, \ie $\prob{I_t > y} = e^{-ty}$ for all $y > 0$.
\end{prop}

\begin{prop}[Mecke equation, {\cite[Th.\,4.5]{LastPenrose18}}]
\label{prel:mecke}
Let $\Pi$ be a Poisson process on $S$ with a $\sigma$-finite mean measure $\mu$, and let $f: \Omega_S \times S^k \to [0, \infty)$ be measurable. Then we have
\begin{multline*}
\EE \sum_{\substack{X_1, \ldots, X_k \\ \text{all distinct}}} f(\Pi\setminus \{X_1, \ldots, X_k\}; X_1, \ldots, X_k) \\[-10pt]
= \int_S \cdots \int_S \EE[f(\Pi; x_1, \ldots, x_k)] \, \mathrm d\mu(x_1)\, \cdots\, \mathrm d\mu(x_k).
\end{multline*}
\end{prop}

\begin{lem}
\label{lem:prob e1c}
We have $\PP[\EC_1^c] \ll \sfrac{1}{\log x}$ for $x\ge2$.
\end{lem}

\begin{proof} We may assume that $x$ is large enough.
We have $\EE[e^{R_x}]\le \EE[e^{\Theta_\infty}] \ll 1$ by \eqref{eq:mgf theta infty}. Therefore, the event $R_x > \log_2 x$ occurs with probability $\ll\sfrac{1}{\log x}$. Using Lemma \ref{lem:scale-invariant-spacing-lemma}, we~find that the probability that there exists $i\in\ZZ$ such that $|S_i-\log x|\le 4r_0$ is $O(1/\log x)$. In~conclusion,
\[
\PP[\EC_1^c] \le \sum_{k\in\{0,1\}} \PP\bgg[4r_0<|S_k-\log x|\le R_x+2r_0,\ R_x\le \log_2x \bgg] +O\bg(\sfrac{1}{\log x}\bg).
\]
\pagebreak[2]
Note that if there exists $k \in \ZZ$ with $W_k>(\log_2x)^{-10}$ and $|S_k-\log x|\le \log_2x + 2r_0$, then $I_{(\log_2x)^{-10}}\ge \log x-\log_2x-2r_0$. So, using Proposition \ref{prop:icexponential}, we~find
\begin{multline*}
\PP\bgg[\exists k: |S_k-\log x|\le \log_2x+2r_0,\ W_k>(\log_2x)^{-10} \bgg] \\
\le \PP\bg[I_{(\log_2 x)^{-10}} > \log x-\log_2x-2r_0\bg]
= \exp\!\bgg(\!-\frac{\log x-\log_2x-2r_0}{(\log_2x)^{10}}\bgg)
\end{multline*}
for $x$ large enough. Consequently,
\begin{multline*}
\PP[\EC_1^c] \le \sum_{k\in\{0,1\}} \PP\bgg[4r_0<|S_k-\log x|\le R_x+2r_0,\ R_x\le \log_2x,\ W_k\le(\log_2x)^{-10} \bgg] \\
+ O\bg(\sfrac{1}{\log x}\bg).
\end{multline*}
Let $m_0$ be the largest integer such that $2^{m_0-1}r_0\le \log_2x$. Hence, if $2r_0<R_x\le\log_2x$, then there exists a unique integer $m\!\in\![2,m_0]$ such that $2^{m-1}r_0\!<\!R_x\!\le\! 2^mr_0$, in~which case the condition $|S_k-\log x|\!\le\! R_x+2r_0$ implies that $|S_k-\log x|\!\le\! 2^mr_0 +2r_0\!\le\! 2^{m+1}r_0$. We~conclude that
\begin{multline*}
\PP[\EC_1^c]
\le \mathop{\sum\sum}_{\substack{k\in\{0,1\} \\ 2\le m\le m_0} }
\PP\bgg[R> 2^{m-1}r_0,\ |S_k-\log x|\le 2^{m+1}r_0,\ W_k \le (\log_2 x)^{-10} \bgg] \\[-10pt]
+O\bg(\sfrac{1}{\log x}\bg).
\end{multline*}
Note that $\one_{R_x > 2^{m-1}r_0}\le \Psfrac{4^{1-m}}{r_0^2} \one_{R_x>2r_0}R_x^2$. In~addition, if $R_x>2r_0$, then we have $\sum_{i\ge2}r(Y_i)^2\le r_0R_x<R_x^2/2$, whence $R_x^2\le 4\sum_{i>j\ge2}r(Y_i)r(Y_j)$.
We conclude that
\begin{multline*}
\PP[\EC_1^c]\le \mathop{\sum\sum}_{\substack{k\in\{0,1\} \\ 2\le m\le m_0} }
\frac{4^{2-m}}{r_0^2} \EE\bggg[ \mathop{\sum\sum}_{i,j :\ i>j>k}r(Y_i)r(Y_j) \cdot \one_{|S_k-\log x|\le 2^{m+1}r_0} \cdot \one_{W_k \le (\log_2x)^{-10}} \bggg] \\[-10pt]
+O\bg(\sfrac{1}{\log x}\bg),
\end{multline*}
Therefore, to complete the proof of the lemma, it is enough to show that
\begin{equation}
\label{eq:E(z)}
E(z) \coloneqq \EE \bggg[\mathop{\sum\sum\sum}_{\substack{i, j, k \in \ZZ:\ i>j>k}}r(Y_i)r(Y_j) \cdot \one_{|S_k-\log x|\le z} \cdot \one_{W_k \le (\log_2x)^{-10}} \bgg] \ll \frac{z}{\log x}
\end{equation}
uniformly for $z\in[0,4\log_2x]$.

For the rest of the proof, we~fix $z\in[0,4\log_2x]$.
Given $t''>t'>t > 0$, let
\[
I_{t,t',t''}\coloneqq \sum_{\substack{(W,Y)\in\RS\\ W\in\RR_{>t}\setminus\{t',t''\}}} Y.
\]
Since the $W_i$'s are almost surely distinct, we~have
\begin{align*}
E(z) \le \EE\bgg[ \mathop{\sum\sum\sum}
r(Y')r(Y'') \cdot \one_{|Y+Y'+Y''+I_{W,W',W''}-\log x|\le z} \cdot \one_{W \le (\log_2x)^{-10}} \bgg].
\end{align*}
with the triple sum being over points $(W,Y),(W',Y'),(W'',Y'')\in\RS$ satisfying \hbox{$W''>W'>W$}. Hence, using the Mecke equation (Proposition \ref{prel:mecke}) and the fact that~$I_{t, t', t''}$ has the same distribution as $I_{t}$, we~find
\[
E(z) \le \mathop{\int\cdots\int}\limits_{\substack{0<w<w'<w'' \\ w\le (\log_2x)^{-10} \\ y,y',y''\ge0 \\ y+y'+y''\le \log x+z }} r(y')r(y'')
\PP\bgg[ \bg|I_w+y+y'+y''-\log x\bg|\le z\bgg]
\frac{\dee w\cdots \dee y''}{ e^{wy+w'y'+w''y''}},
\]
where the integral is sixfold with variables $w,w',w'',y,y',y''$. Proposition \ref{prop:icexponential} implies
\begin{align*}
\PP\bgg[ \bg|I_w+y+y'+y''-\log x\bg|\le z\bgg] &\le e^{-w(\log x-y-y'-y'')} (e^{wz}-e^{-wz}) \\
&\ll wz e^{-w(\log x-y-y'-y'')},
\end{align*}
since $w\le(\log_2x)^{-10}$ and $z\le 4\log_2x$. Consequently,
\[
\begin{split}
E(z) &\ll
\mathop{\int\cdots\int}\limits_{\substack{0<w<w'<w'' \\ y,y',y''\ge0 \\ y+y'+y''\le \log x+z }} r(y')r(y'')
zwe^{-w\log x-(w'-w)(y'+y'')-(w''-w')y''}\dee w\cdots \dee y'' \\
&= \frac{z}{(\log x)^2} \iiiint\limits_{\substack{t,y,y',y''\ge 0 \\ y+y'+y''\le \log x+z }} \frac{r(y')r(y'')}{y''(y'+y'')}
t e^{-t} \dee t\dee y\dee y'\dee y'',
\end{split}
\]
where we made the change of variables $t=w\log x$. Since $\int_0^\infty r(u)/u \dee u\ll 1$ and $\log x+z\ll\log x$, relation \eqref{eq:E(z)} follows. This completes the proof of the lemma.
\end{proof}

\begin{thebibliography}{99}
\bibitem{Arratia98} R. Arratia - “On the central role of scale invariant Poisson processes on (0,∞)”, in Microsurveys in discrete probability (Princeton, NJ, 1997), DIMACS Ser. Discrete Math. Theoret. Comput. Sci., vol. 41, American Mathematical Society, Providence, RI, 1998, p. 21-41
\bibitem{Arratia02} R. Arratia - “On the amount of dependence in the prime factorization of a uniform random integer”, in Contemporary combinatorics, Bolyai Soc. Math. Stud., vol. 10, János Bolyai Math. Soc., Budapest, 2002, p. 29-91
\bibitem{ABT06} R. Arratia, A. D. Barbour \& S. Tavaré - “A tale of three couplings: Poisson-Dirichlet and GEM approximations for random permutations”, Combin. Probab. Comput. 15 (2006) no. 1-2, p. 31-62
\bibitem{BareikisManstavicius07} G. Bareikis \& E. Manstavičius - “On the DDT theorem”, Acta Arith. 126 (2007) no. 2, p. 155-168
\bibitem{Bill72} P. Billingsley - “On the distribution of large prime divisors”, Period. Math. Hungar. 2 (1972), p. 283-289
\bibitem{delaBretecheTenenbaum16} R. de la Bretèche \& G. Tenenbaum - “Sur les processus arithmétiques liés aux diviseurs”, Adv. in Appl. Probab. 48 (2016) no. A, p. 63-76
\bibitem{DDT79} J.-M. Deshouillers, F. Dress \& G. Tenenbaum - “Lois de répartition des diviseurs. I”, Acta Arith. 34 (1979) no. 4, p. 273-285 (loose errata)
\bibitem{DonnellyTavare87} P. Donnelly \& S. Tavaré - “The population genealogy of the infinitely-many neutral alleles model”, J. Math. Biol. 25 (1987) no. 4, p. 381-391
\bibitem{Feng10} S. Feng - The Poisson-Dirichlet distribution and related topics. Models and asymptotic behaviors, Probability and its Appl., Springer, Heidelberg, 2010
\bibitem{JKB} N. L. Johnson, S. Kotz \& N. Balakrishnan - Discrete multivariate distributions, Wiley Series in Probab. and Stat.: Applied Probab. and Stat., John Wiley \& Sons, Inc., New York, 1997
\bibitem{Kingman93} J. F. C. Kingman - Poisson processes, Oxford Studies in Probab., vol. 3, The Clarendon Press, Oxford University Press, New York, 1993
\bibitem{KotzBal00} S. Kotz, N. Balakrishnan \& N. L. Johnson - Continuous multivariate distributions. Vol. 1. Models and applications, Wiley Series in Probab. and Stat.: Applied Probab. and Stat., Wiley-Interscience, New York, 2000
\bibitem{Koukoulo19} D. Koukoulopoulos - The distribution of prime numbers, Graduate Studies in Math., vol. 203, American Mathematical Society, Providence, RI, 2019
\bibitem{LastPenrose18} G. Last \& M. Penrose - Lectures on the Poisson process, Institute of Math. Stat. Textbooks, vol. 7, Cambridge University Press, Cambridge, 2018
\bibitem{Ken22} S.-K. Leung - “Dirichlet law for factorisation of integers, polynomials and permutations”, Math. Proc. Cambridge Philos. Soc. 175 (2023) no. 3, p. 649-676
\bibitem{KenWeb24} S.-K. Leung - “Dirichlet law with arbitrary parameters for factorizations of integers” (2025), Retrieved on 24 August 2025 from https://sites.google.com/view/sunkaileung/miscellaneous
\bibitem{MV07} H. L. Montgomery \& R. C. Vaughan - Multiplicative number theory. I. Classical theory, Cambridge Studies in Advanced Math., vol. 97, Cambridge University Press, Cambridge, 2007
\bibitem{NyandwiSmati13} S. Nyandwi \& A. Smati - “Distribution laws of pairs of divisors”, Integers 13 (2013), article ID A13, 13 pages
\bibitem{Tenenbaum00} G. Tenenbaum - “A rate estimate in Billingsley’s theorem for the size distribution of large prime factors”, Q. J. Math. 51 (2000) no. 3, p. 385-403
\end{thebibliography}
\end{document}
